\documentclass[a4paper, 11pt]{article}

% ── Geometry ──────────────────────────────────────────────────────────
\usepackage[margin=25mm]{geometry}

% ── Maths ─────────────────────────────────────────────────────────────
\usepackage{amsmath, amssymb, bm}
\DeclareMathOperator*{\argmin}{arg\,min}
\DeclareMathOperator{\atantwo}{atan2}
\DeclareMathOperator{\diag}{diag}
\DeclareMathOperator{\trace}{tr}
\newcommand{\R}{\mathbb{R}}
\newcommand{\vect}[1]{\bm{#1}}
\newcommand{\mat}[1]{\bm{#1}}

% ── Typography ────────────────────────────────────────────────────────
\usepackage{parskip}
\usepackage{booktabs}
\usepackage{hyperref}
\hypersetup{colorlinks, linkcolor=blue!60!black, citecolor=blue!60!black, urlcolor=blue!60!black}

% ── Figures ───────────────────────────────────────────────────────────
\usepackage{graphicx}
\usepackage{float}

% ── Code verbatim ─────────────────────────────────────────────────────
\usepackage{listings}
\lstset{basicstyle=\ttfamily\small, breaklines=true, frame=single}

% ── Date ──────────────────────────────────────────────────────────────

% =====================================================================
\title{\textbf{AINF\_SLAMO} \\[4pt]
       \large Implementation--Faithful Mathematical Reference}
\author{Pablo Bustos --- RoboLab}
\date{March 2026}

\begin{document}
\maketitle
\tableofcontents
\newpage

% =====================================================================
\section{Scope}
% =====================================================================

This document catalogues every mathematical equation that is \emph{actually
implemented} in the \texttt{ainf\_slamo} codebase.
The focus is on concrete formulas used in \texttt{C++}, not generic textbook
derivations.
The companion document \texttt{main.tex} provides the theoretical Active
Inference background.

The areas covered are:

\begin{enumerate}[label=(\roman*)]
  \item Frames, state and conventions (Section~\ref{sec:frames}).
  \item Active Inference localiser (Section~\ref{sec:localiser}).
  \item Polygon visibility--graph planner (Section~\ref{sec:planner}).
  \item MPPI/PD trajectory controller (Section~\ref{sec:controller}).
  \item Camera model and synthetic rendering (Section~\ref{sec:camera}).
  \item Object footprints and rotation conventions (Section~\ref{sec:footprints}).
  \item Infrastructure masking (Section~\ref{sec:masking}).
  \item LiDAR processing (Section~\ref{sec:lidar}).
  \item YOLO/segmentation grounding (Section~\ref{sec:grounding}).
  \item Episodic memory (Section~\ref{sec:episodic}).
  \item Navigation completion (Section~\ref{sec:navcompl}).
\end{enumerate}

% =====================================================================
\section{Frames, State and Conventions}\label{sec:frames}
% =====================================================================

All subsystems in \texttt{ainf\_slamo} share a small set of coordinate
frames and conventions.  This section defines the three reference frames
(room, robot, camera), the kinematic model that propagates the robot
state, and the rotation matrix used throughout the codebase.  Consistent
adherence to these definitions ensures that LiDAR, odometry, camera and
planning data can be composed without sign or axis errors.

\subsection{Coordinate frames}

Three frames are used throughout the system:

\begin{description}
  \item[Room (world)] A fixed 2D/3D frame whose $XY$-plane is the floor.
    Polygon vertices and furniture positions live here.
  \item[Robot] Origin at the robot centre, $Y^+$ forward, $X^+$ right.
    LiDAR and odometry are delivered in this frame.
  \item[Camera] Origin at the camera optical centre ($0.92\,$m above the
    robot centre, no lateral offset).
    $X$ lateral, $Y$ forward (depth), $Z$ up.
\end{description}

\subsection{Robot kinematic convention}

The controller and rollout model use $Y^+$ forward, $X^+$ right.
Per-step motion model:
\begin{equation}\label{eq:motion-model}
\begin{aligned}
  x_{t+1} &= x_t + v_t \sin\vartheta_t\,\Delta t, \\
  y_{t+1} &= y_t + v_t \cos\vartheta_t\,\Delta t, \\
  \vartheta_{t+1} &= \vartheta_t + \omega_t\,\Delta t.
\end{aligned}
\end{equation}

\subsection{Localisation state vector}

The room model exposes
\begin{equation}
  \vect{s} = [w,\; l,\; x,\; y,\; \phi]^\top,
\end{equation}
where $w, l$ are room dimensions and $(x, y, \phi)$ is the robot pose in the
room frame.  After initialisation, optimisation acts only on $(x, y, \phi)$.

\subsection{Standard 2D rotation matrix}

Unless stated otherwise every rotation uses
\begin{equation}\label{eq:R2d}
  \mat{R}(\theta)
  = \begin{bmatrix}
      \cos\theta & -\sin\theta \\
      \sin\theta &  \cos\theta
    \end{bmatrix}.
\end{equation}

% =====================================================================
\section{Localiser Mathematics (Active Inference)}\label{sec:localiser}
% =====================================================================

The localiser estimates the robot pose $(x, y, \phi)$ inside a
known room polygon by minimising a variational free energy that
balances LiDAR--wall fit (likelihood) against motion--prediction
priors.  It fuses command and odometry priors via Gaussian precision
weighting, propagates uncertainty through an EKF-style prediction
step, and updates the posterior covariance with a diagonal Fisher
approximation.  Adaptive gradient scaling and early-exit logic keep
the per-frame cost low when the estimate is already trustworthy.

\subsection{LiDAR to room transformation}

Given LiDAR points $\vect{p}_i^{\text{robot}} \in \R^3$,
their $XY$ projection is transformed to the room frame by
\begin{equation}
  \vect{p}_{i,xy}^{\text{room}}
  = \mat{R}(\phi)\,\vect{p}_{i,xy}^{\text{robot}} + \vect{t},
  \qquad
  \vect{t} = [x, y]^\top.
\end{equation}

\subsection{SDF observation model}

For polygon room mode the signed distance is
\begin{equation}\label{eq:sdf}
  d_i = \min_j \bigl\|\vect{p}_i - \Pi_{[a_j, b_j]}(\vect{p}_i)\bigr\|,
\end{equation}
where $\Pi_{[a,b]}(\vect{p})$ is the projection of $\vect{p}$ onto segment
$[a, b]$ with clamped parameter:
\begin{equation}\label{eq:seg-proj}
  t = \operatorname{clamp}\!\Bigl(
    \frac{(\vect{p} - \vect{a}) \cdot (\vect{b} - \vect{a})}
         {\|\vect{b} - \vect{a}\|^2},\; 0,\; 1
  \Bigr),
  \qquad
  \Pi_{[a,b]}(\vect{p}) = \vect{a} + t\,(\vect{b}-\vect{a}).
\end{equation}

\subsection{Variational objective (implemented loss)}

The scalar minimised by the optimiser is
\begin{equation}\label{eq:free-energy}
  \mathcal{F}(\vect{s})
  = \underbrace{
      \frac{1}{2\sigma_{\text{obs}}^2}\;
      \overline{\mathrm{Huber}_\delta(d_i, 0)}
    }_{\text{likelihood term}}
  \;+\;
  \underbrace{
      \frac{1}{2}
      (\vect{s}_{\text{pose}} - \vect{\mu})^\top
      \mat{\Pi}\,
      (\vect{s}_{\text{pose}} - \vect{\mu})
    }_{\text{prior term}},
\end{equation}
with implemented constants:
\begin{center}
\begin{tabular}{ll}
  \toprule
  Symbol & Value \\
  \midrule
  $\sigma_{\text{obs}}$ & $0.05\,$m \\
  $\delta$ (Huber threshold) & $0.15\,$m \\
  $\vect{s}_{\text{pose}}$ & $[x, y, \phi]^\top$ \\
  $(\vect{\mu}, \mat{\Pi})$ & from predicted motion prior \\
  \bottomrule
\end{tabular}
\end{center}

The robust UI metric is
\begin{equation}
  \operatorname{median}_i |d_i|.
\end{equation}

\subsection{Motion integration priors}

Two priors are integrated between LiDAR timestamps:
(i) \emph{command prior} from commanded $(adv_x, adv_z, rot)$ and
(ii) \emph{odometry prior} from feedback $(adv, side, rot)$.
Both integrate piecewise-constant commands over overlap windows:
\begin{equation}\label{eq:motion-int}
  \Delta\vect{p}
  = \sum_k \mat{R}(\theta_k)
    \begin{bmatrix} \Delta x_k^{\text{local}} \\ \Delta y_k^{\text{local}} \end{bmatrix},
  \qquad
  \Delta\theta = \sum_k \Delta\theta_k.
\end{equation}

\subsection{Dual-prior Gaussian fusion}

From command prior $(\vect{\mu}_c, \mat{\Sigma}_c)$ and odometry prior
$(\vect{\mu}_o, \mat{\Sigma}_o)$:
\begin{equation}\label{eq:gaussian-fusion}
\begin{aligned}
  \mat{\Pi}_c &= \mat{\Sigma}_c^{-1}, \qquad
  \mat{\Pi}_o = \mat{\Sigma}_o^{-1}, \\
  \mat{\Pi}_f &= \mat{\Pi}_c + \mat{\Pi}_o, \\
  \vect{\mu}_f &= \mat{\Pi}_f^{-1}\bigl(
    \mat{\Pi}_c\,\vect{\mu}_c + \mat{\Pi}_o\,\vect{\mu}_o
  \bigr).
\end{aligned}
\end{equation}
Angular differences are unwrapped before fusion and the fused angle is
normalised to $[-\pi, \pi]$.

\subsection{EKF-style covariance prediction}

\begin{equation}\label{eq:cov-pred}
  \mat{P}^{-} = \mat{F}\,\mat{P}\,\mat{F}^\top + \mat{Q},
\end{equation}
with Jacobian entries (localised-state case):
\begin{equation}
  F_{x,\theta} = -\Delta y_{\text{loc}}\sin\theta - \Delta x_{\text{loc}}\cos\theta,
  \qquad
  F_{y,\theta} =  \Delta y_{\text{loc}}\cos\theta - \Delta x_{\text{loc}}\sin\theta.
\end{equation}
Noise $\mat{Q}$ is anisotropic in robot-local axes (forward/lateral), then
rotated to global $XY$.  Forward uncertainty is larger than lateral.

\subsection{Covariance update after optimisation}

Posterior precision is
\begin{equation}\label{eq:cov-update}
  \mat{\Lambda}^{+}
  = (\mat{P}^{-})^{-1} + \mat{H}_{\text{like}} + \lambda\,\mat{I},
  \qquad
  \mat{P}^{+} = (\mat{\Lambda}^{+})^{-1},
\end{equation}
where the implemented diagonal Fisher-style approximation is
\begin{equation}
  H_{ii}
  \approx
  \Bigl|\frac{\partial\mathcal{L}_{\text{like}}}{\partial\theta_i}\Bigr|
  \cdot \frac{N}{\mathcal{L}_{\text{like}} + \epsilon}
  + 10^{-4}.
\end{equation}

\subsection{Velocity-adaptive gradient weighting}

During Adam optimisation, gradient components are scaled by
motion-profile weights $w_x, w_y, w_\theta$:
\begin{itemize}
  \item Rotation-dominant: upweight $\theta$, downweight $(x, y)$.
  \item Translation-dominant: upweight the dominant translation axis,
        downweight $\theta$.
\end{itemize}

\subsection{Prediction-based early exit}

Optimisation is skipped when all conditions hold:
\begin{enumerate}
  \item Tracking is stable for a minimum number of steps.
  \item Predicted covariance trace $\trace(\mat{P}^{-})$ is below threshold.
  \item Mean absolute predicted SDF satisfies
    \begin{equation}
      \bar{d} < \tau_{\text{exit}}
      = \sigma_{\text{sdf}} \cdot \text{prediction\_trust\_factor}.
    \end{equation}
\end{enumerate}

\subsection{Innovation and smoothing}

After update, pose is exponentially smoothed:
\begin{equation}
  \hat{\vect{s}}_k
  = \alpha\,\hat{\vect{s}}_{k-1} + (1-\alpha)\,\vect{s}_k,
\end{equation}
with wrapped-angle handling for $\phi$.  Innovation:
\begin{equation}
  \vect{\nu}_k = \vect{s}_k - \vect{\mu}_k^{-},
\end{equation}
with the angular component normalised to $[-\pi, \pi]$.

% =====================================================================
\section{Planner Mathematics (Polygon Visibility Graph)}\label{sec:planner}
% =====================================================================

The global planner constructs a visibility graph over the room polygon
and all obstacle polygons.  Vertices are offset inward (room boundary)
or outward (obstacles) to guarantee collision clearance.  Edges are
added only between mutually visible nodes, verified by a multi-stage
test against room walls and expanded obstacles.  A Dijkstra search on
this weighted graph yields the shortest collision-free path from the
robot to the goal.

\subsection{Point-in-polygon test (ray casting)}

For point $\vect{p}$ and each edge $(v_i, v_j)$, toggle inside/outside
when
\begin{equation}
  \bigl((v_i^y > p_y) \ne (v_j^y > p_y)\bigr)
  \;\land\;
  \Bigl(p_x < \frac{(v_j^x - v_i^x)(p_y - v_i^y)}{v_j^y - v_i^y} + v_i^x\Bigr).
\end{equation}

\subsection{Segment intersection}

Proper intersection uses 2D cross products:
\begin{equation}
  t = \frac{(\vect{b}_1 - \vect{a}_1) \times \vect{d}_2}
           {\vect{d}_1 \times \vect{d}_2},
  \qquad
  u = \frac{(\vect{b}_1 - \vect{a}_1) \times \vect{d}_1}
           {\vect{d}_1 \times \vect{d}_2},
\end{equation}
with strict interior checks $t, u \in (\epsilon, 1-\epsilon)$,
$\epsilon = 10^{-5}$.

\subsection{Inward and outward polygon offsets}

For each vertex, candidates are sampled on a circle of radius $r$ over
$36$ angles.
\begin{itemize}
  \item \textbf{Inward offset}: choose the candidate \emph{inside} the
        polygon that maximises clearance.
  \item \textbf{Outward offset} (obstacles): choose the candidate
        \emph{outside} the polygon that maximises clearance.
\end{itemize}
If no feasible candidate is found, shorter radii are tried.
Clearance score:
\begin{equation}
  c(\vect{p})
  = \min_{e \in \partial\mathcal{P}}
    \operatorname{dist}(\vect{p}, e).
\end{equation}

\subsection{Visibility predicate}

A segment $(a, b)$ is visible iff:
\begin{enumerate}
  \item Both endpoints inside the room polygon.
  \item Endpoints and midpoint not inside any expanded obstacle.
  \item No intersection with the original room boundary.
  \item No intersection with the inner boundary.
  \item No intersection with obstacle edges.
  \item No sampled interior point along the segment lies inside an expanded
        obstacle.
\end{enumerate}

\subsection{Graph construction and shortest path}

Nodes: inner polygon points plus obstacle navigation nodes.
Edge weights are Euclidean:
\begin{equation}
  w_{ij} = \|\vect{p}_i - \vect{p}_j\|_2.
\end{equation}
Start and goal are inserted as temporary nodes with visibility-based
edges.  Shortest path is found by Dijkstra:
\begin{equation}
  \operatorname{dist}(v)
  = \min_{u \to v}\bigl(\operatorname{dist}(u) + w_{uv}\bigr).
\end{equation}
The path is rejected if total length exceeds
\texttt{max\_path\_length}.

% =====================================================================
\section{Controller Mathematics (MPPI + Safety)}\label{sec:controller}
% =====================================================================

The trajectory controller converts the planned path into velocity
commands via Model Predictive Path Integral (MPPI) control.  It rolls
out hundreds of sampled trajectories over a robot-centric Euclidean
Signed Distance Field (ESDF), scores each rollout by a composite cost
that blends goal progress, obstacle proximity, lateral clearance,
control smoothness and an information-theoretic correction, then
weights the samples by a softmax to extract the optimal command.
Adaptive temperature and noise scaling react to the current difficulty
of the environment, while a forward safety guard provides a last-resort
collision avoidance layer.  A PD fallback mode is available for simple
open-space navigation.

\subsection{ESDF construction}

An occupancy grid in robot frame is built from LiDAR plus static obstacle
samples transformed from room frame.
A two-pass chamfer-like propagation approximates the distance transform;
distances are then scaled by grid resolution.

Query: bilinear interpolation.
Gradient: centred finite differences, normalised.

\subsection{Adaptive safety distance near goal}

The effective safety distance interpolates between far and near values:
\begin{equation}\label{eq:dsafe}
  \alpha = e^{-d_{\text{goal}}/\tau},
  \qquad
  d_{\text{safe}}^{\text{eff}}
  = (1-\alpha)\,d_{\text{far}} + \alpha\,d_{\text{near}}.
\end{equation}

\subsection{Obstacle step cost}

Given ESDF value $d$:

\paragraph{Soft term} (inside safety zone):
\begin{equation}
  g_{\text{soft}}
  = \Bigl(\frac{d_{\text{safe}}^{\text{eff}} - d}
               {d_{\text{safe}}^{\text{eff}} - r_{\text{robot}}}\Bigr)^2_+.
\end{equation}

\paragraph{Hard-close term} (inside close margin):
\begin{equation}
  g_{\text{hard}}
  = k_{\text{close}}
    \Bigl(\frac{(r_{\text{robot}} + m_{\text{close}}) - d}
               {m_{\text{close}}}\Bigr)^2_+.
\end{equation}

\paragraph{Final capped cost:}
\begin{equation}
  G_{\text{obs}}^{\text{step}}
  = \min\bigl(
      \lambda_{\text{obs}}(g_{\text{soft}} + g_{\text{hard}}),\;
      G_{\text{obs}}^{\text{cap}}
    \bigr).
\end{equation}

\subsection{MPPI seed generation}

The sampling set includes:
\begin{itemize}
  \item Nominal seed toward carrot.
  \item Six injected exploratory seeds with heading offsets
        $\pm 30°, \pm 60°, \pm 90°$.
  \item Gaussian-perturbed seeds:
    \begin{equation}
      \vect{u}_t^{(k)} = \vect{u}_t^{\text{nom}} + \vect{\epsilon}_t^{(k)},
      \qquad
      \epsilon_{\text{adv}} \sim \mathcal{N}(0, \sigma_{\text{adv}}^2),
      \quad
      \epsilon_{\text{rot}} \sim \mathcal{N}(0, \sigma_{\text{rot}}^2).
    \end{equation}
\end{itemize}

\subsection{Rollout objective}

For each seed the total score is
\begin{equation}\label{eq:rollout-score}
  G = G_{\text{goal}} + G_{\text{obs}} + G_{\text{lat}}
    + G_{\text{smooth}} + G_{\text{vel,mag}} + G_{\text{vel},\Delta}
    + G_{\text{info}}
    + \lambda_{\text{goal}} G_{\text{progress}}
    + \mathbf{1}_{\text{collision}}\,P_{\text{collision}},
\end{equation}
where:
\begin{align}
  G_{\text{goal}}
    &: \text{closest approach to carrot + heading penalty,} \\
  G_{\text{obs}}
    &: \text{discounted sum of per-step obstacle costs,} \\
  G_{\text{lat}}
    &: \text{low side clearance + worsening side trend,} \\
  G_{\text{smooth}}
    &= \lambda_{\text{smooth}}
       (\Delta v_0^2 + \Delta\omega_0^2)
       \;\text{(against previous optimum),} \\
  G_{\text{vel,mag}}
    &= \lambda_v \sum_t (v_t^2 + c_\omega\,\omega_t^2), \\
  G_{\text{vel},\Delta}
    &= \lambda_{\Delta v}
       \sum_t (\Delta v_t^2 + c_\omega\,\Delta\omega_t^2).
\end{align}
Near-term collisions (within a hard horizon) make the seed infeasible;
farther collisions incur a softer penalty.

\subsection{Information-theoretic MPPI correction}

For stochastic seeds:
\begin{equation}
  G_{\text{info}}
  = \lambda \sum_t
    \biggl(
      u_{t,\text{adv}}^{\text{nom}}
      \frac{\epsilon_{t,\text{adv}}}{\sigma_{\text{adv}}^2}
      +
      u_{t,\text{rot}}^{\text{nom}}
      \frac{\epsilon_{t,\text{rot}}}{\sigma_{\text{rot}}^2}
    \biggr).
\end{equation}

\subsection{MPPI weighting, ESS and control extraction}

For non-colliding rollouts:
\begin{equation}
  w_k = \exp\!\Bigl(-\frac{G_k - G_{\min}}{\lambda_{\text{used}}}\Bigr),
  \qquad
  \hat{\vect{u}}_t = \sum_k \bar{w}_k\,\vect{u}_t^{(k)},
  \qquad
  \bar{w}_k = \frac{w_k}{\sum_j w_j}.
\end{equation}
Effective sample size:
\begin{equation}
  \mathrm{ESS}
  = \frac{\bigl(\sum_k w_k\bigr)^2}{\sum_k w_k^2}.
\end{equation}
The command issued averages the first 3 optimised steps and applies
velocity EMA.

\subsection{Dominance-based adaptation}

A dominance signal is computed as
\begin{equation}
  D = p_{\text{free}}\,R_\theta,
  \qquad
  E = 1 - D,
\end{equation}
where $p_{\text{free}}$ is the free-rollout ratio and $R_\theta$ the
steering concentration.  A frontal Safety-Guard gate scales exploration
influence.

Adaptation laws:
\begin{align}
  \lambda &\leftarrow \operatorname{clip}\!\bigl(
    \lambda_0(1 + 2E),\; \lambda_{\min},\; \lambda_{\max}\bigr), \\
  \sigma_{\text{rot}} &\uparrow \text{ with } E, \qquad
  \sigma_{\text{adv}} \downarrow \text{ with } E, \qquad
  T \uparrow \text{ with } E.
\end{align}

\subsection{Safety guard -- forward risk simulation}

A short-horizon forward simulation evaluates risk by querying ESDF along
predicted states.  Trigger conditions:
\begin{itemize}
  \item Hard-threshold violation, or
  \item Enough consecutive soft-threshold violations.
\end{itemize}

If triggered, the controller attempts (in order):
\begin{enumerate}
  \item Backup + turn command.
  \item Backup only.
  \item Scaled forward reduction.
  \item Stop + turn fallback.
\end{enumerate}

\subsection{PD fallback mode}

\begin{equation}
  e_\theta = \atantwo(x_{\text{carrot}}, y_{\text{carrot}}),
  \qquad
  \omega = K_p\,e_\theta + K_d\,\Delta e_\theta,
\end{equation}
\begin{equation}
  v = v_{\max}\,\max(0, \cos e_\theta)^p \cdot
      \text{dist\_factor},
\end{equation}
followed by the same smoothing, Gaussian turning brake and ESDF safety
gate.

% =====================================================================
\section{Camera Model and Synthetic Rendering}\label{sec:camera}
% =====================================================================

The camera subsystem bridges the 3D scene graph and the 2D image
plane.  A calibrated pinhole model (intrinsics derived from the
sensor's horizontal field of view) handles projection and
unprojection.  The synthetic renderer transforms room walls and
furniture into camera coordinates, clips against the near plane,
projects to pixels, and performs ray--plane occlusion tests to
determine visibility.  An attention mode scores each object by
angular proximity to the optical axis, projected size, visibility
and distance, selecting the most salient item for the robot's focus.

\subsection{Camera intrinsics}

The pinhole model is parametrised by a struct
\texttt{CameraIntrinsics} $= (f_x, f_y, c_x, c_y, w, h)$.

Values are computed from the simulated sensor's horizontal field of view
$\theta_h$ at boot time:
\begin{equation}\label{eq:intrinsics}
  f_x = f_y = \frac{w}{2\tan(\theta_h / 2)},
  \qquad
  c_x = \tfrac{w}{2},
  \qquad
  c_y = \tfrac{h}{2}.
\end{equation}
Equality $f_x = f_y$ assumes square pixels.
For the Webots ZED camera: $w = 1280$, $h = 720$,
$\theta_h = 1.9198\,$rad, giving $f_x = f_y \approx 464.8$.

\subsection{World $\to$ robot $\to$ camera transforms}

\paragraph{World to robot.}
\begin{equation}\label{eq:world2robot}
  \vect{p}_{\text{robot}}
  = \mat{R}(-\theta)\,\bigl(\vect{p}_{\text{world}} - \vect{t}\bigr),
  \qquad
  \vect{t} = [x_r, y_r]^\top.
\end{equation}

\paragraph{Robot to camera.}
The camera is mounted at extrinsic offset
$(t_x, t_y, t_z) = (0, 0, 0.92)\,$m (height only):
\begin{equation}\label{eq:robot2cam}
  \vect{p}_{\text{cam}}
  = \begin{bmatrix}
      p_{\text{robot}}^x - t_x \\
      p_{\text{robot}}^y - t_y \\
      z_{\text{world}} - t_z
    \end{bmatrix}.
\end{equation}
Camera axes: $x$ lateral, $y$ forward (depth), $z$ up.

\subsection{Pinhole projection (3D $\to$ 2D)}

A point $\vect{p}_{\text{cam}} = (X, Y, Z)$ with $Y > 0.05\,$m (near plane)
projects to pixel
\begin{equation}\label{eq:projection}
  u = c_x + f_x \,\frac{X}{Y},
  \qquad
  v = c_y - f_y \,\frac{Z}{Y}.
\end{equation}

\subsection{Ray unprojection (2D $\to$ 3D ray)}

Given pixel $(u, v)$ the unit-depth camera ray is
\begin{equation}\label{eq:unproject}
  \vect{d}_{\text{cam}}
  = \biggl(\frac{u - c_x}{f_x},\; 1,\; \frac{c_y - v}{f_y}\biggr).
\end{equation}

\subsection{Near-plane clipping}

For a 3D line segment $(\vect{a}, \vect{b})$, if $a_y \leq 0.05\,$m the
endpoint is interpolated forward:
\begin{equation}
  t = \frac{0.05 - a_y}{b_y - a_y},
  \qquad
  \vect{a}' = \vect{a} + t\,(\vect{b} - \vect{a}).
\end{equation}
The segment is rejected if both endpoints fail the near-plane test.
A margin bounds check (image $\pm 200\,$px) rejects fully off-screen
segments.

\subsection{Wall plane construction}

For each polygon edge from $\vect{a}_0$ to $\vect{b}_0$ (floor level), a
quad is formed with top vertices $\vect{a}_1, \vect{b}_1$ at wall
height~$h$.  The plane normal is
\begin{equation}
  \vect{u} = \vect{b}_0 - \vect{a}_0,
  \qquad
  \vect{v} = \vect{a}_1 - \vect{a}_0,
  \qquad
  \hat{\vect{n}} = \frac{\vect{u} \times \vect{v}}
                        {\|\vect{u} \times \vect{v}\|},
\end{equation}
with plane equation $\hat{\vect{n}} \cdot \vect{p} + d = 0$,
\;$d = -\hat{\vect{n}} \cdot \vect{a}_0$.

\subsection{Ray--plane intersection for occlusion}

A ray from the camera to object centroid $\vect{p}_{\text{cam}}$ hits a
wall plane at
\begin{equation}
  t = \frac{-d}{\hat{\vect{n}} \cdot \vect{p}_{\text{cam}}}.
\end{equation}
The object is visible if $t = \infty$ (no hit) or $t \geq 1 - 10^{-3}$
(wall behind object).

\subsection{Attention mode -- angular proximity scoring}\label{sec:attention}

For each furniture item with centroid $\vect{c}_{\text{cam}}$ in camera
frame, the absolute bearing to the optical axis is
\begin{equation}
  \phi = \bigl|\atantwo(c_{\text{cam}}^x,\; c_{\text{cam}}^y)\bigr|.
\end{equation}
Items with $\phi > 50°$ are rejected.  A composite score selects the best:
\begin{equation}\label{eq:attention-score}
  S = 0.58\,s_{\text{center}}
    + 0.22\,s_{\text{visible}}
    + 0.15\,s_{\text{size}}
    + 0.05\,s_{\text{proximity}},
\end{equation}
where
\begin{align}
  s_{\text{center}}
    &= \operatorname{clamp}\!\bigl(1 - \phi / \phi_{\text{ref}},\; 0,\; 1\bigr),
      \qquad \phi_{\text{ref}} = 20°, \\
  s_{\text{visible}}
    &= \operatorname{clamp}\!\bigl(\text{hits}/6,\; 0,\; 1\bigr), \\
  s_{\text{size}}
    &= \operatorname{clamp}\!\bigl(\ln(A+1)/2,\; 0,\; 1\bigr), \\
  s_{\text{proximity}}
    &= \frac{1}{1 + 0.35\,d},
\end{align}
with $A$ the projected area and $d$ the Euclidean distance.

% =====================================================================
\section{Object Footprints and Rotation Conventions}\label{sec:footprints}
% =====================================================================

Every piece of furniture in the scene graph is represented by a 2D
footprint polygon in a local frame, transformed to world coordinates
via the standard rotation.  Different object types (table, chair, pot)
use different polygon templates.  Two rotation conventions coexist in
the codebase: the standard $R(\theta)$ matrix used by footprints and
renderers, and a non-standard axis convention internal to
\texttt{SceneGraphModel}.  The mapping to Qt3D coordinates adds a
further axis permutation and a $\pi$-offset on the Y-rotation.

\subsection{Local-to-world transform}

A local point $\vect{p}_{\text{loc}}$ is mapped to world by
\begin{equation}\label{eq:local2world}
  \vect{p}_{\text{world}}
  = \vect{t} + \mat{R}(\theta)\,\vect{p}_{\text{loc}},
  \qquad
  \text{i.e.}\quad
  \begin{bmatrix} p_x \\ p_y \end{bmatrix}
  =
  \begin{bmatrix} t_x \\ t_y \end{bmatrix}
  +
  \begin{bmatrix}
    \cos\theta & -\sin\theta \\
    \sin\theta &  \cos\theta
  \end{bmatrix}
  \begin{bmatrix} x_{\text{loc}} \\ y_{\text{loc}} \end{bmatrix}.
\end{equation}
Local $X$ = width, local $Y$ = depth (front direction).
Front in world: $(-\sin\theta, \cos\theta)$.

\subsection{Rectangular footprint}

Four corners at local positions $(\pm w/2,\; \pm d/2)$ transformed
by~\eqref{eq:local2world}.

\subsection{Chair footprint}

An $8$-vertex polygon: a rectangle with a backrest notch at $-Y$ (rear).
Backrest thickness: $b_t = \texttt{backrest\_ratio} \cdot d$.

\subsection{Circular footprint (pot)}

A $16$-gon with radius $r = \max(w, d)/2$:
\begin{equation}
  \vect{v}_i
  = \bigl(r\cos(2\pi i/16),\; r\sin(2\pi i/16)\bigr),
  \qquad i = 0, \ldots, 15.
\end{equation}

\subsection{Rotation conventions}

\paragraph{Standard (footprints, rendering, viewers).}
Uses $\mat{R}(\theta)$ from~\eqref{eq:R2d}.

\paragraph{Non-standard (SceneGraphModel, internal only).}
\begin{equation}
  \vect{y}_{\text{dir}} = (\cos\theta,\; \sin\theta),
  \qquad
  \vect{x}_{\text{dir}} = (\sin\theta,\; -\cos\theta).
\end{equation}
\textbf{Warning:} never use \texttt{SceneGraphModel::local\_to\_world\_polygon()}
for rendering; use \texttt{rc::footprints::make()} instead.

\subsection{Qt3D mapping}

\paragraph{Position.}
Room $(x, y)$ maps to Qt3D as
\begin{equation}
  X_{\text{Qt3D}} = -x, \qquad
  Y_{\text{Qt3D}} = h,  \qquad
  Z_{\text{Qt3D}} = y.
\end{equation}

\paragraph{Y-rotation.}
Qt3D $R_y(\alpha)$ maps $(1,0,0) \to (\cos\alpha, 0, -\sin\alpha)$.
To match the standard-convention width axis
$(\cos\theta, \sin\theta) \to (-\cos\theta, 0, \sin\theta)$ one needs
\begin{equation}\label{eq:qt3d-rot}
  \alpha_{\text{Qt3D}} = \pi + \theta.
\end{equation}

% =====================================================================
\section{Infrastructure Masking}\label{sec:masking}
% =====================================================================

Infrastructure masking segments each camera pixel into
\emph{infrastructure} (floor, walls, ceiling) or \emph{foreground}
(furniture, people).  For every pixel a ray is cast from the camera
through the pinhole model and intersected analytically with the room's
floor plane, ceiling plane and wall segments.  The nearest valid hit
gives the infrastructure depth; pixels whose measured depth is close to
or beyond this value are classified as infrastructure and can be
suppressed in the display.

For each camera pixel $(u, v)$ a ray
$\vect{d}_{\text{cam}}$ is constructed via~\eqref{eq:unproject} and
intersected with floor, ceiling and wall planes.
The infrastructure depth is the minimum valid hit:
\begin{equation}
  d_{\text{infra}} = \min\bigl(t_{\text{floor}},\; t_{\text{ceil}},\;
                               t_{\text{walls}}\bigr).
\end{equation}

\subsection{Floor intersection}

\begin{equation}
  t_{\text{floor}}
  = \frac{-z_{\text{cam}}}{d_{\text{cam}}^z},
  \qquad
  \text{valid if } |d_{\text{cam}}^z| > 10^{-6}
  \;\text{and}\;
  \vect{p}_{xy} \in \text{polygon}.
\end{equation}

\subsection{Ceiling intersection}

\begin{equation}
  t_{\text{ceil}}
  = \frac{h_{\text{wall}} - z_{\text{cam}}}{d_{\text{cam}}^z}.
\end{equation}

\subsection{Wall intersection (2D ray--segment)}

For each wall edge $(\vect{a}, \vect{b})$, using 2D cross products:
\begin{equation}
  \operatorname{cross2}(\vect{u}, \vect{v}) = u_x v_y - u_y v_x,
\end{equation}
\begin{equation}
  t = \frac{\operatorname{cross2}(\vect{q} - \vect{p},\; \vect{b} - \vect{a})}
           {\operatorname{cross2}(\vect{d}_{\text{ray}},\; \vect{b} - \vect{a})},
  \qquad
  s_u = \frac{\operatorname{cross2}(\vect{q} - \vect{p},\; \vect{d}_{\text{ray}})}
             {\operatorname{cross2}(\vect{d}_{\text{ray}},\; \vect{b} - \vect{a})}.
\end{equation}
Valid if $t > 0.05\,$m, $s_u \in [0, 1]$ and
$0 \leq z_{\text{hit}} \leq h_{\text{wall}}$.

% =====================================================================
\section{LiDAR Processing}\label{sec:lidar}
% =====================================================================

Raw LiDAR point clouds are filtered, decimated and transformed before
they reach the localiser or controller.  Height-band filters separate
wall-level returns (used for room fitting) from low-level returns (used
for obstacle detection).  A stride-based decimation keeps the point
count bounded.  Finally, a rigid SE(2) transform maps robot-frame
points into the room frame for use by the SDF observation model.

\subsection{Height filtering}

\begin{center}
\begin{tabular}{lll}
  \toprule
  Sensor & Height band & Purpose \\
  \midrule
  HELIOS (high) & $z \in [1.2, 2.2]\,$m & Wall detection \\
  BPEARL (low)  & $z > 0.1\,$m          & Obstacle/people \\
  \bottomrule
\end{tabular}
\end{center}
Maximum range: $100\,$m.

\subsection{Decimation}

\begin{equation}
  \text{stride}
  = \max\!\bigl(1,\;
    \lfloor\text{num\_points}\,/\,\text{pool\_size}\rfloor\bigr).
\end{equation}

\subsection{LiDAR $\to$ room transform (SE(2))}

\begin{equation}
  \vect{p}_{\text{room}}^{xy}
  = \mat{R}(\phi)\,\vect{p}_{\text{robot}}^{xy} + \vect{t},
  \qquad
  \vect{t} = [x, y]^\top.
\end{equation}

\subsection{LiDAR visualisation in Qt3D}

World coordinates:
\begin{equation}
  w_x = \cos\theta\, p_x - \sin\theta\, p_y + x_r,
  \qquad
  w_z = \sin\theta\, p_x + \cos\theta\, p_y + y_r.
\end{equation}
Mapped to Qt3D: $(-w_x,\; p_z,\; w_z)$.

% =====================================================================
\section{YOLO / Segmentation Grounding}\label{sec:grounding}
% =====================================================================

Detections from an external YOLO-based image segmentation service are
grounded into the 3D scene graph.  Each detection mask is back-projected
through the depth buffer to obtain a 3D point cloud, which is
transformed from camera to robot to room frame.  Ground-plane points
are clipped, and a PCA over the remaining cloud yields the object
centroid and principal axis for pose estimation.  This section also
covers the depth-buffer format conversions that handle the various
encoding schemes.

\subsection{3D point transformation pipeline}

\begin{enumerate}
  \item \textbf{Camera $\to$ robot}:
    $\vect{p}_{\text{robot}}
    = \vect{p}_{\text{cam}} + (\Delta t_x, \Delta t_y, \Delta t_z)$.
  \item \textbf{Robot $\to$ room}:
    $\vect{p}_{\text{room}} = \mat{R}(\theta)\,\vect{p}_{\text{robot}} + \vect{t}$.
  \item \textbf{Z-lift}: $z' = z_{\text{robot}} + 0.05\,$m (keep above floor).
  \item Output: $\vect{p}_{3\text{D}} = (p_{\text{room}}^x,\; p_{\text{room}}^y,\; z')$.
\end{enumerate}

\subsection{Ground clipping}

\begin{equation}
  z_{\text{floor}} = \min_i z_i + 0.06\,\text{m},
  \qquad
  \text{keep } \{p \mid p_z > z_{\text{floor}}\}.
\end{equation}
If the filtered set has $< 8$ points, the original cloud is kept.

\subsection{Centroid and covariance}

\begin{equation}
  \vect{m} = \frac{1}{N}\sum_{i=1}^N \vect{p}_i,
  \qquad
  \mat{\Sigma}
  = \frac{1}{N}\sum_{i=1}^N
    (\vect{p}_i - \vect{m})(\vect{p}_i - \vect{m})^\top.
\end{equation}
Major axis = eigenvector of $\mat{\Sigma}$ with largest eigenvalue.

\subsection{Depth buffer extraction}

\begin{center}
\begin{tabular}{ll}
  \toprule
  Format & Conversion \\
  \midrule
  float32 & direct copy, apply scale \\
  uint16  & mm $\to$ m \\
  uint8   & raw bytes, scale to metres \\
  \bottomrule
\end{tabular}
\end{center}
Invalid values (NaN, Inf) are mapped to $0$.

% =====================================================================
\section{Episodic Memory}\label{sec:episodic}
% =====================================================================

The episodic memory records navigation episodes as feature vectors
summarising trajectory statistics, speed profiles, control efficiency
and computational load.  At query time, the $k$ most similar past
episodes are retrieved via a weighted Euclidean distance in feature
space.  Supporting geometric primitives (point-to-segment and
point-to-polygon distance) are reused from the planner and shared with
the masking subsystem.

\subsection{Point-to-segment distance}

\begin{equation}
  \vect{ab} = \vect{b} - \vect{a},
  \qquad
  t = \operatorname{clamp}\!\Bigl(
    \frac{(\vect{p} - \vect{a}) \cdot \vect{ab}}
         {\max(10^{-9},\;\|\vect{ab}\|^2)},\; 0,\; 1
  \Bigr),
  \qquad
  d = \bigl\|\vect{p} - (\vect{a} + t\,\vect{ab})\bigr\|.
\end{equation}

\subsection{Point-to-polygon distance}

\begin{equation}
  d_{\text{poly}}
  = \min_i \operatorname{dist}(\vect{p},\; \text{edge}_i).
\end{equation}

\subsection{Episode similarity ($k$-NN retrieval)}

Weighted Euclidean distance over feature vectors:
\begin{equation}
  d(\vect{f}_1, \vect{f}_2)
  = \sqrt{\sum_i w_i\,(f_{1,i} - f_{2,i})^2}.
\end{equation}
The $k = 5$ nearest episodes are retrieved.

\subsection{Stored episode metrics}

\begin{center}
\begin{tabular}{ll}
  \toprule
  Category & Fields \\
  \midrule
  Trajectory    & \texttt{distance\_traveled, path\_efficiency} \\
  Speed         & \texttt{mean\_speed, p95\_speed} \\
  Rotation      & \texttt{mean\_rot, p95\_rot} \\
  Efficiency    & \texttt{mean\_ess\_ratio, p05\_ess\_ratio} \\
  Performance   & \texttt{mean\_cpu\_pct, p95\_mppi\_ms} \\
  \bottomrule
\end{tabular}
\end{center}

% =====================================================================
\section{Navigation Completion for \texttt{goto\_object}}\label{sec:navcompl}
% =====================================================================

Once the trajectory controller declares the path complete, a final
heading-alignment phase rotates the robot to face the target object.
This phase is decoupled from trajectory tracking and uses a simple
yaw-error criterion with timeout and anti-oscillation guards.

After path completion in object-navigation mode, a dedicated final heading
phase (decoupled from trajectory tracking) steers toward the target.
The yaw error in robot frame is
\begin{equation}
  e_\theta = \atantwo(x_{\text{obj}}^{\text{robot}},\;
                      y_{\text{obj}}^{\text{robot}}).
\end{equation}
Exit conditions combine angle tolerance, distance checks, timeout cycles
and anti-oscillation guards.

% =====================================================================
\section{Room Geometry}\label{sec:room-geom}
% =====================================================================

The room is represented as a closed 2D polygon whose vertices define
the floor boundary.  This polygon serves multiple roles: it is the
target shape for localisation SDF queries, the outer boundary for the
visibility-graph planner, the source of wall planes for infrastructure
masking and synthetic rendering, and the container for
point-in-polygon membership tests.  Wall meshes for 3D visualisation
are extruded from the polygon edges to a specified ceiling height.

\subsection{Polygon representation}

The room boundary is stored as a closed polygon
$\{\vect{v}_i\}_{i=0}^{N-1} \subset \R^2$.

\subsection{Signed distance field (polygon)}

See Equation~\eqref{eq:sdf} with clamped segment projection~\eqref{eq:seg-proj}.

\subsection{Wall mesh (3D)}

Per edge $(x_1, y_1) \to (x_2, y_2)$:
\begin{equation}
  \text{Quad vertices:}\quad
  \begin{cases}
    v_0 = (x_1, y_1, 0), & v_1 = (x_1, y_1, h), \\
    v_2 = (x_2, y_2, 0), & v_3 = (x_2, y_2, h).
  \end{cases}
\end{equation}
Triangulated as $(v_0, v_2, v_1)$ and $(v_1, v_2, v_3)$.
Total: $N$ edges $\to$ $6N$ triangle indices.

% =====================================================================
\section{Symbol Glossary}\label{sec:glossary}
% =====================================================================

\begin{center}
\begin{tabular}{lp{10cm}}
  \toprule
  Symbol & Meaning \\
  \midrule
  $x, y, \theta$           & Robot pose in controller/localisation space \\
  $v, \omega$               & Advance speed and rotational speed \\
  $\Delta t$                & Control/update time step \\
  $\vect{s} = [w,l,x,y,\phi]^\top$ & Room-localisation state \\
  $d_i$                     & SDF distance residual for LiDAR point $i$ \\
  $\sigma_{\text{obs}}, \delta$ & Observation noise scale and Huber threshold \\
  $\vect{\mu}, \mat{\Pi}$   & Prior mean and precision matrix \\
  $\mat{P}^{-}, \mat{P}^{+}$ & Predicted and posterior covariance \\
  $\mat{F}, \mat{Q}$        & Motion Jacobian and process-noise covariance \\
  $\vect{\nu}$              & Innovation residual \\
  $f_x, f_y, c_x, c_y$     & Camera intrinsic parameters \\
  $\theta_h$                & Camera horizontal field of view \\
  $d_{\text{safe}}^{\text{eff}}$ & Effective safety distance \\
  $G$                       & Total trajectory score for a rollout \\
  $\lambda$                 & MPPI temperature (context-dependent) \\
  $\mathrm{ESS}$            & Effective sample size of rollout weights \\
  $D, E$                    & Dominance and exploration signals (adaptive MPPI) \\
  $\mat{R}(\theta)$         & Standard 2D rotation matrix \\
  \bottomrule
\end{tabular}
\end{center}

% =====================================================================
\section{Equation to Implementation Mapping}\label{sec:mapping}
% =====================================================================

\subsection{Localiser}

\begin{center}
\small
\begin{tabular}{p{7cm}p{7cm}}
  \toprule
  Equation & Implementation \\
  \midrule
  Variational loss $\mathcal{F}$ \eqref{eq:free-energy}
    & \texttt{loss\_sdf\_mse(...)} + \texttt{m.prior\_loss()} \\
  Prior quadratic term
    & \texttt{Model::prior\_loss()} with \texttt{prediction\_precision\_matrix} \\
  Dual-prior fusion \eqref{eq:gaussian-fusion}
    & \texttt{RoomConceptAI::fuse\_priors(...)} \\
  Covariance prediction \eqref{eq:cov-pred}
    & \texttt{predict\_step(...)} \\
  Covariance update \eqref{eq:cov-update}
    & covariance update block in \texttt{update(...)} \\
  Early-exit threshold
    & \texttt{sigma\_sdf} $\times$ \texttt{prediction\_trust\_factor} \\
  \bottomrule
\end{tabular}
\end{center}

\subsection{Planner}

\begin{center}
\small
\begin{tabular}{p{7cm}p{7cm}}
  \toprule
  Equation & Implementation \\
  \midrule
  Point-in-polygon   & \texttt{point\_in\_polygon(...)} \\
  Segment intersection & \texttt{segments\_intersect\_proper(...)} \\
  Inward/outward offsets & \texttt{offset\_polygon\_inward/outward(...)} \\
  Visibility          & \texttt{is\_visible(...)} \\
  Shortest path       & \texttt{dijkstra(...)}, \texttt{plan(...)} \\
  \bottomrule
\end{tabular}
\end{center}

\subsection{Controller}

\begin{center}
\small
\begin{tabular}{p{7cm}p{7cm}}
  \toprule
  Equation & Implementation \\
  \midrule
  Obstacle step cost  & \texttt{obstacle\_step\_cost(...)} \\
  Repulsion strength  & \texttt{obstacle\_repulsion\_strength(...)} \\
  Rollout score \eqref{eq:rollout-score}
                      & \texttt{simulate\_and\_score(...)} \\
  Info correction     & \texttt{simulate\_and\_score} with \texttt{use\_info\_correction} \\
  MPPI control + ESS  & \texttt{compute(...)}, \texttt{compute\_ess(...)} \\
  Dominance adaptation & \texttt{adapt\_from\_dominance(...)} \\
  Safety gate         & safety gate block in \texttt{compute(...)}, \texttt{compute\_pd(...)} \\
  \bottomrule
\end{tabular}
\end{center}

\subsection{Camera and rendering}

\begin{center}
\small
\begin{tabular}{p{7cm}p{7cm}}
  \toprule
  Equation & Implementation \\
  \midrule
  Intrinsics \eqref{eq:intrinsics}
    & \texttt{CameraViewer} intrinsics callback \\
  Projection \eqref{eq:projection}
    & \texttt{SyntheticCameraRenderer::project\_point()} \\
  Unprojection \eqref{eq:unproject}
    & ray construction in \texttt{CameraViewer::grab\_frame()} \\
  Near-plane clip
    & \texttt{SyntheticCameraRenderer::clip\_and\_project()} \\
  Attention score \eqref{eq:attention-score}
    & \texttt{pick\_attention\_target()} \\
  Qt3D rotation \eqref{eq:qt3d-rot}
    & object transform in \texttt{Viewer3D} \\
  \bottomrule
\end{tabular}
\end{center}

% =====================================================================
\appendix
\section{Build and Run}
% =====================================================================

\begin{lstlisting}
cmake --fresh -S . -B build_make -G "Unix Makefiles"
make -C build_make -j20
bin/ainf_slamo etc/config
\end{lstlisting}

\end{document}
